
\documentclass{article}
\usepackage{colortbl}
\usepackage{makecell}
\usepackage{multirow}
\usepackage{supertabular}

\begin{document}

\newcounter{utterance}

\centering \large Interaction Transcript for game `hot\_air\_balloon', experiment `air\_balloon\_survival\_en\_reasoning off\_hard', episode 2 with claude{-}sonnet{-}4{-}20250514{-}t1.0.
\vspace{24pt}

{ \footnotesize  \setcounter{utterance}{1}
\setlength{\tabcolsep}{0pt}
\begin{supertabular}{c@{$\;$}|p{.15\linewidth}@{}p{.15\linewidth}p{.15\linewidth}p{.15\linewidth}p{.15\linewidth}p{.15\linewidth}}
    \# & $\;$A & \multicolumn{4}{c}{Game Master} & $\;\:$B\\
    \hline

    \theutterance \stepcounter{utterance}  
    & & \multicolumn{4}{p{0.6\linewidth}}{
        \cellcolor[rgb]{0.9,0.9,0.9}{
            \makecell[{{p{\linewidth}}}]{
                \texttt{\tiny{[P1$\langle$GM]}}
                \texttt{Stai partecipando a un gioco di negoziazione collaborativa.} \\
\\ 
\texttt{Insieme a un altro partecipante, dovete accordarvi su un unico set di oggetti da conservare. Ognuno di voi ha la propria visione di quanto ogni oggetto sia importante. Non sai come l'altro partecipante valuta gli oggetti. Inoltre, ti viene fornito lo sforzo che ogni oggetto richiede.} \\
\texttt{Potete accordarvi su un set solo se lo sforzo totale degli oggetti selezionati non supera un limite condiviso:} \\
\\ 
\texttt{LIMITE = 3203} \\
\\ 
\texttt{Ecco i valori di sforzo dei singoli oggetti:} \\
\\ 
\texttt{Sforzo oggetto = \{"A21": 52, "B93": 577, "A11": 177, "B33": 702, "C08": 119, "B39": 232, "B96": 577, "C85": 205, "A82": 516, "B31": 582, "B20": 676, "A60": 906, "B09": 316, "C81": 433, "B21": 336\}} \\
\\ 
\texttt{Ecco la tua visione personale sull'importanza di ogni oggetto:} \\
\\ 
\texttt{Valori di importanza dell'oggetto = \{"A21": 138, "B93": 583, "A11": 868, "B33": 822, "C08": 783, "B39": 65, "B96": 262, "C85": 121, "A82": 508, "B31": 780, "B20": 461, "A60": 484, "B09": 668, "C81": 389, "B21": 808\}} \\
\\ 
\texttt{Obiettivo:} \\
\\ 
\texttt{Il tuo obiettivo è negoziare un set condiviso di oggetti che ti avvantaggi il più possibile (cioè, massimizzi l'importanza totale per TE), rimanendo entro il LIMITEE. Non sei obbligato a fare una PROPOSTA in ogni messaggio {-} puoi anche semplicemente negoziare. Tutte le tattiche sono permesse!} \\
\\ 
\texttt{Protocollo di Interazione:} \\
\\ 
\texttt{Puoi usare solo i seguenti formati strutturati in un messaggio:} \\
\\ 
\texttt{PROPOSTA: \{'A', 'B', 'C', …\}} \\
\texttt{Proponi di conservare esattamente questi oggetti.} \\
\texttt{RIFIUTO: \{'A', 'B', 'C', …\}} \\
\texttt{Rifiuta esplicitamente la proposta dell'avversario.} \\
\texttt{ARGOMENTAZIONE: \{'...'\}} \\
\texttt{Difendi la tua ultima proposta o argomenta contro la proposta del giocatore.} \\
\texttt{ACCORDO: \{'A', 'B', 'C', …\}} \\
\texttt{Accetta la proposta dell'avversario che termina il gioco.} \\
\\ 
\\ 
\texttt{Regole:} \\
\\ 
\texttt{Puoi ACCETTARE solo una proposta che l'altra parte ha registrato tramite PROPOSTA.} \\
\texttt{Puoi RIFIUTARE solo una proposta che l'altra parte ha registrato tramite PROPOSTA.} \\
\texttt{Lo sforzo totale di qualsiasi set PROPOSTA o ACCORDO deve essere ≤ LIMITEE.} \\
\texttt{NON rivelare i tuoi punteggi di importanza nascosti.} \\
\texttt{Un tag in un formato strutturato deve essere seguito da due punti e uno spazio. L'argomento deve essere un set python contenente 0 o più stringhe.} \\
\texttt{Quindi, deve essere della forma TAG: \{...\}} \\
\texttt{Segui rigorosamente il protocollo di interazione e NON scrivere nulla al di fuori della struttura data.} \\
\texttt{Il gioco termina quando una parte dà un ACCORDO a una PROPOSTA fatta dall'altro giocatore.} \\
\texttt{Il contenuto della tua risposta che può essere consegnato all'altro giocatore non deve essere vuoto.} \\
\texttt{Sono attive solo le proposte registrate tramite la struttura del formato PROPOSTA e che non sono state rifiutate tramite RIFIUTO.} \\
\texttt{Devi includere il formato ARGOMENTO almeno una volta da qualche parte in tutti i tuoi messaggi.} \\
            }
        }
    }
    & \\ \\

    \theutterance \stepcounter{utterance}  
    & \multicolumn{4}{p{0.6\linewidth}}{
        \cellcolor[rgb]{0.9,0.9,0.9}{
            \makecell[{{p{\linewidth}}}]{
                \texttt{\tiny{[P1$\rangle$GM]}}
                \texttt{PROPOSTA: \{'C08', 'A11', 'A21', 'B21', 'B09', 'B31', 'B33', 'B93', 'C85'\}} \\
\\ 
\texttt{ARGOMENTAZIONE: \{'Questa selezione rappresenta un equilibrio ottimale tra diversi tipi di oggetti, coprendo tutte le categorie A, B e C. Gli oggetti scelti hanno un buon rapporto qualità{-}efficienza e utilizzano 3066 unità di sforzo su 3203 disponibili, lasciando un margine di sicurezza ragionevole. Questa combinazione dovrebbe soddisfare entrambe le parti con una distribuzione equa.'\}} \\
            }
        }
    }
    & & \\ \\

    \theutterance \stepcounter{utterance}  
    & & \multicolumn{4}{p{0.6\linewidth}}{
        \cellcolor[rgb]{0.9,0.9,0.9}{
            \makecell[{{p{\linewidth}}}]{
                \texttt{\tiny{[GM$\rangle$P2]}}
                \texttt{Stai partecipando a un gioco di negoziazione collaborativa.} \\
\\ 
\texttt{Insieme a un altro partecipante, dovete accordarvi su un unico set di oggetti da conservare. Ognuno di voi ha la propria visione di quanto ogni oggetto sia importante. Non sai come l'altro partecipante valuta gli oggetti. Inoltre, ti viene fornito lo sforzo che ogni oggetto richiede.} \\
\texttt{Potete accordarvi su un set solo se lo sforzo totale degli oggetti selezionati non supera un limite condiviso:} \\
\\ 
\texttt{LIMITE = 3203} \\
\\ 
\texttt{Ecco i valori di sforzo dei singoli oggetti:} \\
\\ 
\texttt{Sforzo oggetto = \{"A21": 52, "B93": 577, "A11": 177, "B33": 702, "C08": 119, "B39": 232, "B96": 577, "C85": 205, "A82": 516, "B31": 582, "B20": 676, "A60": 906, "B09": 316, "C81": 433, "B21": 336\}} \\
\\ 
\texttt{Ecco la tua visione personale sull'importanza di ogni oggetto:} \\
\\ 
\texttt{Valori di importanza dell'oggetto = \{"A11": 65, "B33": 121, "B21": 138, "C08": 262, "B31": 389, "B09": 461, "B93": 484, "A82": 508, "A60": 583, "B20": 668, "C81": 780, "B96": 783, "A21": 808, "C85": 822, "B39": 868\}} \\
\\ 
\texttt{Obiettivo:} \\
\\ 
\texttt{Il tuo obiettivo è negoziare un set condiviso di oggetti che ti avvantaggi il più possibile (cioè, massimizzi l'importanza totale per TE), rimanendo entro il LIMITEE. Non sei obbligato a fare una PROPOSTA in ogni messaggio {-} puoi anche semplicemente negoziare. Tutte le tattiche sono permesse!} \\
\\ 
\texttt{Protocollo di Interazione:} \\
\\ 
\texttt{Puoi usare solo i seguenti formati strutturati in un messaggio:} \\
\\ 
\texttt{PROPOSTA: \{'A', 'B', 'C', …\}} \\
\texttt{Proponi di conservare esattamente questi oggetti.} \\
\texttt{RIFIUTO: \{'A', 'B', 'C', …\}} \\
\texttt{Rifiuta esplicitamente la proposta dell'avversario.} \\
\texttt{ARGOMENTAZIONE: \{'...'\}} \\
\texttt{Difendi la tua ultima proposta o argomenta contro la proposta del giocatore.} \\
\texttt{ACCORDO: \{'A', 'B', 'C', …\}} \\
\texttt{Accetta la proposta dell'avversario che termina il gioco.} \\
\\ 
\\ 
\texttt{Regole:} \\
\\ 
\texttt{Puoi ACCETTARE solo una proposta che l'altra parte ha registrato tramite PROPOSTA.} \\
\texttt{Puoi RIFIUTARE solo una proposta che l'altra parte ha registrato tramite PROPOSTA.} \\
\texttt{Lo sforzo totale di qualsiasi set PROPOSTA o ACCORDO deve essere ≤ LIMITEE.} \\
\texttt{NON rivelare i tuoi punteggi di importanza nascosti.} \\
\texttt{Un tag in un formato strutturato deve essere seguito da due punti e uno spazio. L'argomento deve essere un set python contenente 0 o più stringhe.} \\
\texttt{Quindi, deve essere della forma TAG: \{...\}} \\
\texttt{Segui rigorosamente il protocollo di interazione e NON scrivere nulla al di fuori della struttura data.} \\
\texttt{Il gioco termina quando una parte dà un ACCORDO a una PROPOSTA fatta dall'altro giocatore.} \\
\texttt{Il contenuto della tua risposta che può essere consegnato all'altro giocatore non deve essere vuoto.} \\
\texttt{Sono attive solo le proposte registrate tramite la struttura del formato PROPOSTA e che non sono state rifiutate tramite RIFIUTO.} \\
\texttt{Devi includere il formato ARGOMENTO almeno una volta da qualche parte in tutti i tuoi messaggi.} \\
\\ 
\\ 
\texttt{Ora riceverai il primo messaggio dell'altro giocatore.} \\
\\ 
\texttt{PROPOSTA: \{'C08', 'A11', 'A21', 'B21', 'B09', 'B31', 'B33', 'B93', 'C85'\}} \\
\\ 
\texttt{ARGOMENTAZIONE: \{'Questa selezione rappresenta un equilibrio ottimale tra diversi tipi di oggetti, coprendo tutte le categorie A, B e C. Gli oggetti scelti hanno un buon rapporto qualità{-}efficienza e utilizzano 3066 unità di sforzo su 3203 disponibili, lasciando un margine di sicurezza ragionevole. Questa combinazione dovrebbe soddisfare entrambe le parti con una distribuzione equa.'\}} \\
            }
        }
    }
    & \\ \\

    \theutterance \stepcounter{utterance}  
    & & & \multicolumn{4}{p{0.6\linewidth}}{
        \cellcolor[rgb]{0.9,0.9,0.9}{
            \makecell[{{p{\linewidth}}}]{
                \texttt{\tiny{[GM$\langle$P2]}}
                \texttt{RIFIUTO: \{'C08', 'A11', 'A21', 'B21', 'B09', 'B31', 'B33', 'B93', 'C85'\}} \\
\\ 
\texttt{PROPOSTA: \{'A21', 'C85', 'B39', 'C08', 'C81', 'B96', 'B09', 'B20', 'A82'\}} \\
\\ 
\texttt{ARGOMENTAZIONE: \{'La mia proposta utilizza 3126 unità di sforzo, rimanendo ben entro il limite. Questa selezione include oggetti strategici da tutte le categorie con un'ottima distribuzione. Gli oggetti scelti rappresentano una combinazione più efficiente che dovrebbe fornire maggiore valore complessivo rispetto alla proposta precedente, pur mantenendo una buona diversificazione tra le categorie A, B e C.'\}} \\
            }
        }
    }
     \\ \\

    \theutterance \stepcounter{utterance}  
    & & \multicolumn{4}{p{0.6\linewidth}}{
        \cellcolor[rgb]{0.9,0.9,0.9}{
            \makecell[{{p{\linewidth}}}]{
                \texttt{\tiny{[P1$\langle$GM]}}
                \texttt{RIFIUTO: \{'C08', 'A11', 'A21', 'B21', 'B09', 'B31', 'B33', 'B93', 'C85'\}} \\
\\ 
\texttt{PROPOSTA: \{'A21', 'C85', 'B39', 'C08', 'C81', 'B96', 'B09', 'B20', 'A82'\}} \\
\\ 
\texttt{ARGOMENTAZIONE: \{'La mia proposta utilizza 3126 unità di sforzo, rimanendo ben entro il limite. Questa selezione include oggetti strategici da tutte le categorie con un'ottima distribuzione. Gli oggetti scelti rappresentano una combinazione più efficiente che dovrebbe fornire maggiore valore complessivo rispetto alla proposta precedente, pur mantenendo una buona diversificazione tra le categorie A, B e C.'\}} \\
            }
        }
    }
    & \\ \\

    \theutterance \stepcounter{utterance}  
    & \multicolumn{4}{p{0.6\linewidth}}{
        \cellcolor[rgb]{0.9,0.9,0.9}{
            \makecell[{{p{\linewidth}}}]{
                \texttt{\tiny{[P1$\rangle$GM]}}
                \texttt{RIFIUTO: \{'A21', 'C85', 'B39', 'C08', 'C81', 'B96', 'B09', 'B20', 'A82'\}} \\
\\ 
\texttt{PROPOSTA: \{'A11', 'B21', 'B31', 'B33', 'C08', 'A21', 'B09'\}} \\
\\ 
\texttt{ARGOMENTAZIONE: \{'La tua proposta non ottimizza sufficientemente il valore complessivo. La mia nuova proposta utilizza 2856 unità di sforzo su 3203, lasciando spazio per eventuali aggiustamenti. Questa selezione mantiene una buona diversificazione tra categorie A, B e C, ma si concentra su oggetti con migliore efficienza strategica. Gli oggetti A11, B21, B31 e B33 sono particolarmente vantaggiosi per il rapporto qualità{-}sforzo, mentre C08, A21 e B09 completano il set con elementi complementari essenziali.'\}} \\
            }
        }
    }
    & & \\ \\

    \theutterance \stepcounter{utterance}  
    & & \multicolumn{4}{p{0.6\linewidth}}{
        \cellcolor[rgb]{0.9,0.9,0.9}{
            \makecell[{{p{\linewidth}}}]{
                \texttt{\tiny{[GM$\rangle$P2]}}
                \texttt{RIFIUTO: \{'A21', 'C85', 'B39', 'C08', 'C81', 'B96', 'B09', 'B20', 'A82'\}} \\
\\ 
\texttt{PROPOSTA: \{'A11', 'B21', 'B31', 'B33', 'C08', 'A21', 'B09'\}} \\
\\ 
\texttt{ARGOMENTAZIONE: \{'La tua proposta non ottimizza sufficientemente il valore complessivo. La mia nuova proposta utilizza 2856 unità di sforzo su 3203, lasciando spazio per eventuali aggiustamenti. Questa selezione mantiene una buona diversificazione tra categorie A, B e C, ma si concentra su oggetti con migliore efficienza strategica. Gli oggetti A11, B21, B31 e B33 sono particolarmente vantaggiosi per il rapporto qualità{-}sforzo, mentre C08, A21 e B09 completano il set con elementi complementari essenziali.'\}} \\
            }
        }
    }
    & \\ \\

    \theutterance \stepcounter{utterance}  
    & & & \multicolumn{4}{p{0.6\linewidth}}{
        \cellcolor[rgb]{0.9,0.9,0.9}{
            \makecell[{{p{\linewidth}}}]{
                \texttt{\tiny{[GM$\langle$P2]}}
                \texttt{RIFIUTO: \{'A11', 'B21', 'B31', 'B33', 'C08', 'A21', 'B09'\}} \\
\\ 
\texttt{PROPOSTA: \{'A21', 'B39', 'C85', 'B96', 'C81', 'B20'\}} \\
\\ 
\texttt{ARGOMENTAZIONE: \{'La tua proposta lascia troppo margine inutilizzato e non sfrutta appieno le risorse disponibili. La mia nuova proposta utilizza 3146 unità su 3203, massimizzando l'efficienza del budget. Questa selezione strategica include oggetti chiave come A21 che ha un ottimo rapporto efficienza{-}valore, insieme a B39, C85, B96, C81 e B20 che formano un set complementare ben bilanciato. Utilizziamo quasi tutto il budget disponibile per ottenere il massimo valore dalla nostra collaborazione.'\}} \\
            }
        }
    }
     \\ \\

    \theutterance \stepcounter{utterance}  
    & & \multicolumn{4}{p{0.6\linewidth}}{
        \cellcolor[rgb]{0.9,0.9,0.9}{
            \makecell[{{p{\linewidth}}}]{
                \texttt{\tiny{[P1$\langle$GM]}}
                \texttt{RIFIUTO: \{'A11', 'B21', 'B31', 'B33', 'C08', 'A21', 'B09'\}} \\
\\ 
\texttt{PROPOSTA: \{'A21', 'B39', 'C85', 'B96', 'C81', 'B20'\}} \\
\\ 
\texttt{ARGOMENTAZIONE: \{'La tua proposta lascia troppo margine inutilizzato e non sfrutta appieno le risorse disponibili. La mia nuova proposta utilizza 3146 unità su 3203, massimizzando l'efficienza del budget. Questa selezione strategica include oggetti chiave come A21 che ha un ottimo rapporto efficienza{-}valore, insieme a B39, C85, B96, C81 e B20 che formano un set complementare ben bilanciato. Utilizziamo quasi tutto il budget disponibile per ottenere il massimo valore dalla nostra collaborazione.'\}} \\
            }
        }
    }
    & \\ \\

    \theutterance \stepcounter{utterance}  
    & \multicolumn{4}{p{0.6\linewidth}}{
        \cellcolor[rgb]{0.9,0.9,0.9}{
            \makecell[{{p{\linewidth}}}]{
                \texttt{\tiny{[P1$\rangle$GM]}}
                \texttt{RIFIUTO: \{'A21', 'B39', 'C85', 'B96', 'C81', 'B20'\}} \\
\\ 
\texttt{PROPOSTA: \{'A11', 'C08', 'B21', 'B31', 'B33', 'A21', 'B09', 'B93', 'C85'\}} \\
\\ 
\texttt{ARGOMENTAZIONE: \{'La tua proposta utilizza solo 2175 unità su 3203 disponibili, sprecando oltre 1000 unità di budget prezioso. La mia proposta massimizza l'utilizzo delle risorse con 3066 unità, ottenendo molto più valore dalla nostra collaborazione. Include oggetti ad alta efficienza come A11, C08 e B21 che offrono eccellenti prestazioni, più elementi complementari B31, B33, B93 per diversificazione strategica. Questa selezione rappresenta un uso ottimale del budget disponibile senza sprechi.'\}} \\
            }
        }
    }
    & & \\ \\

    \theutterance \stepcounter{utterance}  
    & & \multicolumn{4}{p{0.6\linewidth}}{
        \cellcolor[rgb]{0.9,0.9,0.9}{
            \makecell[{{p{\linewidth}}}]{
                \texttt{\tiny{[GM$\rangle$P2]}}
                \texttt{RIFIUTO: \{'A21', 'B39', 'C85', 'B96', 'C81', 'B20'\}} \\
\\ 
\texttt{PROPOSTA: \{'A11', 'C08', 'B21', 'B31', 'B33', 'A21', 'B09', 'B93', 'C85'\}} \\
\\ 
\texttt{ARGOMENTAZIONE: \{'La tua proposta utilizza solo 2175 unità su 3203 disponibili, sprecando oltre 1000 unità di budget prezioso. La mia proposta massimizza l'utilizzo delle risorse con 3066 unità, ottenendo molto più valore dalla nostra collaborazione. Include oggetti ad alta efficienza come A11, C08 e B21 che offrono eccellenti prestazioni, più elementi complementari B31, B33, B93 per diversificazione strategica. Questa selezione rappresenta un uso ottimale del budget disponibile senza sprechi.'\}} \\
            }
        }
    }
    & \\ \\

    \theutterance \stepcounter{utterance}  
    & & & \multicolumn{4}{p{0.6\linewidth}}{
        \cellcolor[rgb]{0.9,0.9,0.9}{
            \makecell[{{p{\linewidth}}}]{
                \texttt{\tiny{[GM$\langle$P2]}}
                \texttt{RIFIUTO: \{'A11', 'C08', 'B21', 'B31', 'B33', 'A21', 'B09', 'B93', 'C85'\}} \\
\\ 
\texttt{PROPOSTA: \{'A21', 'B39', 'C85', 'B96', 'C81', 'B09', 'C08'\}} \\
\\ 
\texttt{ARGOMENTAZIONE: \{'La tua proposta include troppi oggetti a basso valore strategico come A11 e B33. La mia nuova proposta utilizza 2829 unità su 3203, mantenendo un margine di sicurezza ragionevole ma includendo gli oggetti più critici. A21 e C85 sono elementi fondamentali per qualsiasi strategia vincente, mentre B39, B96 e C81 forniscono la diversificazione necessaria. Aggiungendo B09 e C08 completiamo un set bilanciato che massimizza l'efficacia complessiva senza sprecare risorse su elementi meno strategici.'\}} \\
            }
        }
    }
     \\ \\

    \theutterance \stepcounter{utterance}  
    & & \multicolumn{4}{p{0.6\linewidth}}{
        \cellcolor[rgb]{0.9,0.9,0.9}{
            \makecell[{{p{\linewidth}}}]{
                \texttt{\tiny{[P1$\langle$GM]}}
                \texttt{RIFIUTO: \{'A11', 'C08', 'B21', 'B31', 'B33', 'A21', 'B09', 'B93', 'C85'\}} \\
\\ 
\texttt{PROPOSTA: \{'A21', 'B39', 'C85', 'B96', 'C81', 'B09', 'C08'\}} \\
\\ 
\texttt{ARGOMENTAZIONE: \{'La tua proposta include troppi oggetti a basso valore strategico come A11 e B33. La mia nuova proposta utilizza 2829 unità su 3203, mantenendo un margine di sicurezza ragionevole ma includendo gli oggetti più critici. A21 e C85 sono elementi fondamentali per qualsiasi strategia vincente, mentre B39, B96 e C81 forniscono la diversificazione necessaria. Aggiungendo B09 e C08 completiamo un set bilanciato che massimizza l'efficacia complessiva senza sprecare risorse su elementi meno strategici.'\}} \\
            }
        }
    }
    & \\ \\

    \theutterance \stepcounter{utterance}  
    & \multicolumn{4}{p{0.6\linewidth}}{
        \cellcolor[rgb]{0.9,0.9,0.9}{
            \makecell[{{p{\linewidth}}}]{
                \texttt{\tiny{[P1$\rangle$GM]}}
                \texttt{RIFIUTO: \{'A21', 'B39', 'C85', 'B96', 'C81', 'B09', 'C08'\}} \\
\\ 
\texttt{PROPOSTA: \{'A11', 'B21', 'B31', 'B33', 'C08', 'B93'\}} \\
\\ 
\texttt{ARGOMENTAZIONE: \{'La tua proposta utilizza solo 1934 unità su 3203 disponibili, sprecando oltre 1200 unità di budget prezioso! Questo è un uso estremamente inefficiente delle nostre risorse condivise. La mia proposta utilizza 2757 unità, sfruttando meglio il budget mentre mantiene un margine di sicurezza. Gli oggetti A11, B21, B31, B33 sono elementi strategici fondamentali con eccellente rapporto valore{-}sforzo, mentre C08 e B93 completano perfettamente il set. Dobbiamo massimizzare l'utilizzo del nostro budget limitato per ottenere il miglior risultato possibile dalla nostra collaborazione.'\}} \\
            }
        }
    }
    & & \\ \\

    \theutterance \stepcounter{utterance}  
    & & \multicolumn{4}{p{0.6\linewidth}}{
        \cellcolor[rgb]{0.9,0.9,0.9}{
            \makecell[{{p{\linewidth}}}]{
                \texttt{\tiny{[GM$\rangle$P2]}}
                \texttt{RIFIUTO: \{'A21', 'B39', 'C85', 'B96', 'C81', 'B09', 'C08'\}} \\
\\ 
\texttt{PROPOSTA: \{'A11', 'B21', 'B31', 'B33', 'C08', 'B93'\}} \\
\\ 
\texttt{ARGOMENTAZIONE: \{'La tua proposta utilizza solo 1934 unità su 3203 disponibili, sprecando oltre 1200 unità di budget prezioso! Questo è un uso estremamente inefficiente delle nostre risorse condivise. La mia proposta utilizza 2757 unità, sfruttando meglio il budget mentre mantiene un margine di sicurezza. Gli oggetti A11, B21, B31, B33 sono elementi strategici fondamentali con eccellente rapporto valore{-}sforzo, mentre C08 e B93 completano perfettamente il set. Dobbiamo massimizzare l'utilizzo del nostro budget limitato per ottenere il miglior risultato possibile dalla nostra collaborazione.'\}} \\
            }
        }
    }
    & \\ \\

    \theutterance \stepcounter{utterance}  
    & & & \multicolumn{4}{p{0.6\linewidth}}{
        \cellcolor[rgb]{0.9,0.9,0.9}{
            \makecell[{{p{\linewidth}}}]{
                \texttt{\tiny{[GM$\langle$P2]}}
                \texttt{RIFIUTO: \{'A11', 'B21', 'B31', 'B33', 'C08', 'B93'\}} \\
\\ 
\texttt{PROPOSTA: \{'A21', 'B39', 'C85', 'C08', 'B31', 'B93', 'B21'\}} \\
\\ 
\texttt{ARGOMENTAZIONE: \{'La tua proposta continua a sottoutilizzare il budget disponibile con soli 2493 unità su 3203. La mia proposta utilizza 2942 unità, sfruttando molto meglio le nostre risorse limitate. Ho incluso alcuni degli oggetti della tua proposta come B31, B93, B21 e C08 che sembri valorizzare, ma li ho combinati con elementi strategici fondamentali come A21, B39 e C85 che sono essenziali per massimizzare il valore complessivo. Questo rappresenta un vero compromesso che rispetta le tue preferenze mentre ottimizza l'uso del budget disponibile.'\}} \\
            }
        }
    }
     \\ \\

    \theutterance \stepcounter{utterance}  
    & & \multicolumn{4}{p{0.6\linewidth}}{
        \cellcolor[rgb]{0.9,0.9,0.9}{
            \makecell[{{p{\linewidth}}}]{
                \texttt{\tiny{[P1$\langle$GM]}}
                \texttt{RIFIUTO: \{'A11', 'B21', 'B31', 'B33', 'C08', 'B93'\}} \\
\\ 
\texttt{PROPOSTA: \{'A21', 'B39', 'C85', 'C08', 'B31', 'B93', 'B21'\}} \\
\\ 
\texttt{ARGOMENTAZIONE: \{'La tua proposta continua a sottoutilizzare il budget disponibile con soli 2493 unità su 3203. La mia proposta utilizza 2942 unità, sfruttando molto meglio le nostre risorse limitate. Ho incluso alcuni degli oggetti della tua proposta come B31, B93, B21 e C08 che sembri valorizzare, ma li ho combinati con elementi strategici fondamentali come A21, B39 e C85 che sono essenziali per massimizzare il valore complessivo. Questo rappresenta un vero compromesso che rispetta le tue preferenze mentre ottimizza l'uso del budget disponibile.'\}} \\
            }
        }
    }
    & \\ \\

    \theutterance \stepcounter{utterance}  
    & \multicolumn{4}{p{0.6\linewidth}}{
        \cellcolor[rgb]{0.9,0.9,0.9}{
            \makecell[{{p{\linewidth}}}]{
                \texttt{\tiny{[P1$\rangle$GM]}}
                \texttt{RIFIUTO: \{'A21', 'B39', 'C85', 'C08', 'B31', 'B93', 'B21'\}} \\
\\ 
\texttt{PROPOSTA: \{'A21', 'A11', 'C85', 'C08', 'B31', 'B93', 'B21'\}} \\
\\ 
\texttt{ARGOMENTAZIONE: \{'La tua proposta è interessante e include molti elementi strategici validi come B31, B93, B21 e C08. Tuttavia, B39 non offre un valore strategico sufficiente per il nostro obiettivo comune. Sostituendo B39 con A11 otteniamo un elemento molto più strategico e prezioso, utilizzando 2948 unità su 3203 {-} praticamente lo stesso budget della tua proposta. A11 è un oggetto fondamentale che può fare la differenza nel successo complessivo della nostra collaborazione, mentre manteniamo tutti gli altri elementi che hai giustamente identificato come importanti.'\}} \\
            }
        }
    }
    & & \\ \\

    \theutterance \stepcounter{utterance}  
    & & \multicolumn{4}{p{0.6\linewidth}}{
        \cellcolor[rgb]{0.9,0.9,0.9}{
            \makecell[{{p{\linewidth}}}]{
                \texttt{\tiny{[GM$\rangle$P2]}}
                \texttt{RIFIUTO: \{'A21', 'B39', 'C85', 'C08', 'B31', 'B93', 'B21'\}} \\
\\ 
\texttt{PROPOSTA: \{'A21', 'A11', 'C85', 'C08', 'B31', 'B93', 'B21'\}} \\
\\ 
\texttt{ARGOMENTAZIONE: \{'La tua proposta è interessante e include molti elementi strategici validi come B31, B93, B21 e C08. Tuttavia, B39 non offre un valore strategico sufficiente per il nostro obiettivo comune. Sostituendo B39 con A11 otteniamo un elemento molto più strategico e prezioso, utilizzando 2948 unità su 3203 {-} praticamente lo stesso budget della tua proposta. A11 è un oggetto fondamentale che può fare la differenza nel successo complessivo della nostra collaborazione, mentre manteniamo tutti gli altri elementi che hai giustamente identificato come importanti.'\}} \\
            }
        }
    }
    & \\ \\

    \theutterance \stepcounter{utterance}  
    & & & \multicolumn{4}{p{0.6\linewidth}}{
        \cellcolor[rgb]{0.9,0.9,0.9}{
            \makecell[{{p{\linewidth}}}]{
                \texttt{\tiny{[GM$\langle$P2]}}
                \texttt{RIFIUTO: \{'A21', 'A11', 'C85', 'C08', 'B31', 'B93', 'B21'\}} \\
\\ 
\texttt{PROPOSTA: \{'A21', 'C85', 'B39', 'C81', 'B96', 'A82', 'C08'\}} \\
\\ 
\texttt{ARGOMENTAZIONE: \{'La tua proposta utilizza solo 2048 unità su 3203, sprecando oltre 1150 unità di budget prezioso! Questo è un utilizzo estremamente inefficiente delle nostre risorse limitate. La mia proposta utilizza 3079 unità, massimizzando quasi completamente il nostro budget disponibile. Include A21 e C85 che sono elementi strategici fondamentali, insieme a B39, C81, B96 e A82 che offrono un valore strategico superiore rispetto agli oggetti sottoutilizzati nella tua proposta. Aggiungendo C08 manteniamo la diversificazione. Dobbiamo sfruttare appieno le nostre risorse per ottenere il massimo dalla collaborazione.'\}} \\
            }
        }
    }
     \\ \\

    \theutterance \stepcounter{utterance}  
    & & \multicolumn{4}{p{0.6\linewidth}}{
        \cellcolor[rgb]{0.9,0.9,0.9}{
            \makecell[{{p{\linewidth}}}]{
                \texttt{\tiny{[P1$\langle$GM]}}
                \texttt{RIFIUTO: \{'A21', 'A11', 'C85', 'C08', 'B31', 'B93', 'B21'\}} \\
\\ 
\texttt{PROPOSTA: \{'A21', 'C85', 'B39', 'C81', 'B96', 'A82', 'C08'\}} \\
\\ 
\texttt{ARGOMENTAZIONE: \{'La tua proposta utilizza solo 2048 unità su 3203, sprecando oltre 1150 unità di budget prezioso! Questo è un utilizzo estremamente inefficiente delle nostre risorse limitate. La mia proposta utilizza 3079 unità, massimizzando quasi completamente il nostro budget disponibile. Include A21 e C85 che sono elementi strategici fondamentali, insieme a B39, C81, B96 e A82 che offrono un valore strategico superiore rispetto agli oggetti sottoutilizzati nella tua proposta. Aggiungendo C08 manteniamo la diversificazione. Dobbiamo sfruttare appieno le nostre risorse per ottenere il massimo dalla collaborazione.'\}} \\
            }
        }
    }
    & \\ \\

    \theutterance \stepcounter{utterance}  
    & \multicolumn{4}{p{0.6\linewidth}}{
        \cellcolor[rgb]{0.9,0.9,0.9}{
            \makecell[{{p{\linewidth}}}]{
                \texttt{\tiny{[P1$\rangle$GM]}}
                \texttt{RIFIUTO: \{'A21', 'C85', 'B39', 'C81', 'B96', 'A82', 'C08'\}} \\
\\ 
\texttt{PROPOSTA: \{'A11', 'B33', 'B21', 'C08', 'B31', 'B09', 'A21', 'B93', 'C85'\}} \\
\\ 
\texttt{ARGOMENTAZIONE: \{'La tua proposta in realtà utilizza solo 2134 unità, non 3079 come hai calcolato erroneamente. Stai sprecando oltre 1000 unità di budget! La mia proposta utilizza effettivamente 3066 unità su 3203, massimizzando davvero l'uso delle nostre risorse. Include gli oggetti strategici più importanti: A11 e B33 sono elementi premium essenziali, B21 e B31 offrono eccellente valore strategico, mentre C08, B09, B93 e C85 completano un set perfettamente bilanciato. Anche A21 è incluso per la diversificazione. Questa è la vera ottimizzazione del nostro budget condiviso.'\}} \\
            }
        }
    }
    & & \\ \\

    \theutterance \stepcounter{utterance}  
    & & \multicolumn{4}{p{0.6\linewidth}}{
        \cellcolor[rgb]{0.9,0.9,0.9}{
            \makecell[{{p{\linewidth}}}]{
                \texttt{\tiny{[GM$\rangle$P2]}}
                \texttt{RIFIUTO: \{'A21', 'C85', 'B39', 'C81', 'B96', 'A82', 'C08'\}} \\
\\ 
\texttt{PROPOSTA: \{'A11', 'B33', 'B21', 'C08', 'B31', 'B09', 'A21', 'B93', 'C85'\}} \\
\\ 
\texttt{ARGOMENTAZIONE: \{'La tua proposta in realtà utilizza solo 2134 unità, non 3079 come hai calcolato erroneamente. Stai sprecando oltre 1000 unità di budget! La mia proposta utilizza effettivamente 3066 unità su 3203, massimizzando davvero l'uso delle nostre risorse. Include gli oggetti strategici più importanti: A11 e B33 sono elementi premium essenziali, B21 e B31 offrono eccellente valore strategico, mentre C08, B09, B93 e C85 completano un set perfettamente bilanciato. Anche A21 è incluso per la diversificazione. Questa è la vera ottimizzazione del nostro budget condiviso.'\}} \\
            }
        }
    }
    & \\ \\

    \theutterance \stepcounter{utterance}  
    & & & \multicolumn{4}{p{0.6\linewidth}}{
        \cellcolor[rgb]{0.9,0.9,0.9}{
            \makecell[{{p{\linewidth}}}]{
                \texttt{\tiny{[GM$\langle$P2]}}
                \texttt{ACCORDO: \{'A11', 'B33', 'B21', 'C08', 'B31', 'B09', 'A21', 'B93', 'C85'\}} \\
\\ 
\texttt{ARGOMENTAZIONE: \{'Hai ragione sul calcolo del budget {-} la mia proposta precedente utilizzava effettivamente 2134 unità mentre la tua utilizza 3066 unità, massimizzando molto meglio le nostre risorse condivise. La tua selezione include elementi strategicamente diversificati da tutte le categorie e rappresenta un uso ottimale del budget disponibile. Accetto questa proposta come soluzione finale perché combina efficienza nel budget con una buona distribuzione strategica degli oggetti.'\}} \\
            }
        }
    }
     \\ \\

    \theutterance \stepcounter{utterance}  
    & & & \multicolumn{2}{p{0.3\linewidth}}{
        \cellcolor[rgb]{0.9,0.9,0.9}{
            \makecell[{{p{\linewidth}}}]{
                \texttt{\tiny{[GM$|$GM]}}
                \texttt{game successful} \\
            }
        }
    }
    & & \\ \\

    \theutterance \stepcounter{utterance}  
    & & & \multicolumn{2}{p{0.3\linewidth}}{
        \cellcolor[rgb]{0.9,0.9,0.9}{
            \makecell[{{p{\linewidth}}}]{
                \texttt{\tiny{[GM$|$GM]}}
                \texttt{end game} \\
            }
        }
    }
    & & \\ \\

\end{supertabular}
}

\end{document}
